\section{Conclusion}

The article contains a self-contained presentation of Hofmann's LFPL.
%
The mechanized soundness proof constructs, for each LFPL expression,
an explicit polynomials that bound the evaluation cost with respect to
a big-step cost semantics.
%
The mechanized completeness proof shows how LFPL can simulate
polynomial-time Turing machines, only relying on linear lists and
functions to encode the polynomial-sized tape in a non-size-increasing
way.
%
Both proofs are based on existing ideas but contain significant
simplifications and improvements.

Our hope is that this work contributes to making LFPL's metatheory
more accessible and inspiring LFPL-based innovations in areas like
automatic resource bound analysis~\cite{HoffmJ22}, quantitative type
theory~\cite{atkeyPolynomialTimeDependent2024}, and memory management
in functional languages~\cite{lorenzenFP2FullyInPlace2023}.

%%% Local Variables:
%%% mode: LaTeX
%%% TeX-master: "../paper"
%%% End:
